%% Cover Letter -- New Submission to Digital Signal Processing
%% Fisher--Rao-Guided Generative Steganography:
%% Geodesic Embedding and Detection Analysis
%% Ekodeck et al.

\documentclass[12pt,a4paper]{article}

\usepackage[utf8]{inputenc}
\usepackage[T1]{fontenc}
\usepackage[margin=2.5cm]{geometry}
\usepackage{microtype}
\usepackage{parskip}
\usepackage{hyperref}
\usepackage{xcolor}

\hypersetup{colorlinks=true, urlcolor=blue, linkcolor=blue}

\pagestyle{empty}

\begin{document}

\raggedright

\textbf{St\'ephane Ga\"el R. Ekodeck}\\
University of Yaounde I -- LIRIMA/GRIMCAPE\\
UMMISCO, Sorbonne Universites, Paris\\
\href{mailto:stephane-gael.ekodeck@facsciences-uy1.cm}
     {stephane-gael.ekodeck@facsciences-uy1.cm}

\bigskip
\today

\bigskip

\textit{Digital Signal Processing}\\
Elsevier

\bigskip

\textbf{Re: Submission of manuscript --
``Fisher--Rao-Guided Generative Steganography: Geodesic Embedding and Detection Analysis''}

\bigskip

Dear Editor,

We are pleased to submit the manuscript entitled
``Fisher--Rao-Guided Generative Steganography: Geodesic Embedding and
Detection Analysis'' for consideration in \textit{Digital Signal Processing}.

The manuscript treats generative and coverless steganography as a statistical
signal-generation and detection problem. Its central question is whether a
message-bearing generator should be analyzed only through the endpoint
divergence between cover and stego distributions, or whether the embedding
trajectory through the generator parameter space also matters for accumulated
detectability. We show that the trajectory does matter: in the small-step
regime, the sequential Kullback--Leibler cost of an embedding path is governed
by Fisher path energy, and Fisher--Rao geodesics provide locally optimal
embedding trajectories.

We believe the paper fits the scope of \textit{Digital Signal Processing}
because it contributes to statistical signal processing, detection theory,
information security, and mathematical methods with direct application to
signal-processing models. Fisher--Rao geometry is not presented as an end in
itself; it is used as the intrinsic sensitivity metric for controlling and
analyzing low-detectability generative steganographic embedding.

The main contributions are:

\begin{itemize}
  \item a detection-aware model of generative steganographic embedding as
        controlled motion on a statistical signal manifold;
  \item a sequential-KL path-cost formulation showing why endpoint divergence
        alone can miss accumulated embedding detectability;
  \item a geodesic embedding result proving local optimality of Fisher--Rao
        paths in the small-step regime;
  \item a Natural-Gradient Steganography algorithm with convergence analysis
        under geodesic smoothness and strong-convexity assumptions;
  \item curvature-based robustness diagnostics for embedding-path sensitivity;
  \item a Neyman--Pearson detection interpretation connecting the steganalyzer's
        statistic to the cover-to-stego displacement in exponential-family
        signal models.
\end{itemize}

The manuscript is methodological and theoretical. It is not submitted as a
complete operational image-steganography pipeline; rather, it provides a
signal-processing framework for designing and analyzing embedding paths in
generative statistical models. Numerical experiments on Gaussian and
Gaussian-mixture signal models illustrate the predicted KL reduction,
natural-gradient speedups, and curvature-driven sensitivity effects.

The manuscript has not been published elsewhere and is not under consideration
by any other journal. All authors have approved this submission. The authors
declare no competing interests. AI-assisted language tools were used to improve
readability; all scientific content was reviewed and edited by the authors, who
take full responsibility for the manuscript.

\bigskip
Respectfully yours,

\medskip
\textbf{St\'ephane Ga\"el R. Ekodeck}\\
Corresponding author\\
\textit{on behalf of all co-authors:}
P.R. Mama Assandje,
S.A. Eb\'el\'e,
S.W. Mossebo Tcheunteu,
A.U. Ewane,
D.B.M. Fotso,
R. Ndoundam.

\end{document}
